
@@M0@@ et la
filtration par le degré sur le deuxième membre résulte aisément de (V 6.6.4) et du fait que les @@M1@@ définissent une \emph{théorie de Chern} de @@M2@@ au sens de (V 6.13), lorsqu'on considère @@M3@@ comme une @@M4@@-algèbre augmentée sur l'anneau binômial @@M5@@). On en déduit des applications
@@M6@@
par la formule
@@M7@@
comparer (XIV 5.1 et 5.2). Ceci posé, la commutativité des diagrammes (7.13.2) et (7.13.4) implique trivialement celle des diagrammes analogues
@@M8@@
@@M9@@
\textbf{Remarque 7.13.11.} La commutativité des deux derniers diagrammes, qu'on peut exprimer en disant que la spécialisation des cycles algébriques est compatible avec la formation des classes de cohomologie @@M10@@-adiques associées, n'est pas très satisfaisante sous la forme énoncée, puisqu'on a été obligé de remplacer les coefficients @@M11@@ par @@M12@@, c'est-à-dire de négliger la torsion. Pour voir un énoncé qui ne néglige pas la torsion, il semble qu'on
